%%%%%
%
\section{Comptage des imagettes de la sonde PMS-2D-C}

\itemstyle{Op\'erateur~:} {\sf Compte\_img\_2D}
\index{Compte\_img\_2D}

%%%%%OBJET%%%%
\itemstyle{Objet~:}
Cet op�rateur calcule le nombre d'imagettes comprises dans les donn�es images
d'une sonde PMS-2D-C ou 2D-P.

%%%%%ARGUMENTS%%%%
\bigskip
%%%%%ENTREES%%%%
\begin{tabarg}
\multicolumn{3}{l}{\titretab {Entr\'ees}}\\
\hline
$Time_{reg}$ & Vecteur, {\tt Float\_64} & Temps avec fr\'equence de 1 Hertz ($s$)\\
$Image_{2D}$ & Matrice X1024, {\tt Int\_32} & Image\\
$Time_{asyn}$ & Vecteur, {\tt Float\_64} & Temps de la variable d'entr\'ee Image ($s$)\\
$TAS_{2D}$ & Vecteur, {\tt Int\_16} & True Air Speed ($m/s$)\\
$Overload_{2D}$ & Vecteur, {\tt Int\_16} & Nombre d'overload\\
$Code\_Hexa\_sonde$ & Coefficient, {\tt String} & Code hexadecimal de la sonde\\
$Masque\_sonde$ & Coefficient, {\tt String} & Masque hexadecimal des donn\'ees images\\
$Bits\_temps$ & Coefficient, {\tt Int\_16} & Nombre de bits de codage du temps dans les s\'eparateurs\\
$Nbval_{as}$ & Coefficient, {\tt Int\_32} & Nombre d'\'el\'ements dans le vecteur de temps ascynchrone\\
\end{tabarg}

\medskip
%%%%%SORTIES%%%%
\begin{tabarg}
\multicolumn{3}{l}{\titretab {Sorties}}\\
\hline
$Compte_{2D}$ & Vecteur, {\tt Int\_32} & Compte total des imagettes 2D\\
$Compte\_vide_{2D}$ & Vecteur, {\tt Int\_32} & Compte des imagettes 2D vides\\
$Compte\_art_{2D}$ & Vecteur, {\tt Int\_32} & Compte des imagettes 2D consid�r�es comme des artifacts �lectroniques\\
\end{tabarg}


%%%%%FORMULE%%%%
\itemstyle{Formule}

%%% Compte_img_2D


%%%%%COMMENTAIRE%%%%
\itemstyle{Commentaire~:}
Le nombre d'imagettes est calcul� � partir du calcul du nombre de s�parateurs
inclus dans les donn�es images. Les r�gles de d�termination des s�parateurs est
moins stricte que dans la d�finition des s�parateurs des donn�es images de la
sonde 2D-C ou 2D-P. Habituellement, les s�parateurs sont constitu�s de la s�quence
suivante (pixel 0 = diode obscurcie): une ou plusieurs ligne(s) de pixels � 1; 
une ligne avec l'information du temps interne cod� sur 24 ou 32 bits (si le
codage est sur 24 bits, les 8 derniers bits contiennent le code
d'identification de la sonde: dans le cas de la sonde du CAM, le code est
101010); une ligne de pixels � 0 sur 24 bits avec les 8
derniers bits contenant le code d'identification de la sonde. Dans notre cas,
nous avons omis la v�rification de la pr�sence d'une ou de plusieurs ligne(s)
de pixels � 1. Cela ne provoque pas d'erreurs dans la d�tection des
s�parateurs et permet de ne pas en omettre.

%%%%%AUTEUR%%%%
\itemstyle{Auteur~:}
Christian Pag\' CNRM/GMEI/MNP

%%%%%REFERENCE%%%%
\itemstyle{R\'ef\'erences~:}
